\documentclass[10pt,a4paper]{amsart}
\usepackage{amsmath,amssymb,amsthm,graphicx}
\usepackage{verbatim}
\newenvironment{reference}{\comment}{\endcomment}
%\newenvironment{reference}{}{}
\newenvironment{slogan}{\comment}{\endcomment}
\newenvironment{history}{\comment}{\endcomment}
% For commutative diagrams we use Xy-pic
\usepackage[all]{xy}

% We use 2cell for 2-commutative diagrams.
\xyoption{2cell}
% Theorem environments.
%
\theoremstyle{plain}
\newtheorem{theorem}[subsection]{Theorem}
\newtheorem{proposition}[subsection]{Proposition}
\newtheorem{lemma}[subsection]{Lemma}

\theoremstyle{definition}
\newtheorem{definition}[subsection]{Definition}
\newtheorem{example}[subsection]{Example}
\newtheorem{exercise}[subsection]{Exercise}
\newtheorem{situation}[subsection]{Situation}

\theoremstyle{remark}
\newtheorem{remark}[subsection]{Remark}
\newtheorem{remarks}[subsection]{Remarks}

\numberwithin{equation}{subsection}

% Macros
%
\def\lim{\mathop{\mathrm{lim}}\nolimits}
\def\colim{\mathop{\mathrm{colim}}\nolimits}
\def\Spec{\mathop{\mathrm{Spec}}}
\def\Hom{\mathop{\mathrm{Hom}}\nolimits}
\def\Ext{\mathop{\mathrm{Ext}}\nolimits}
\def\SheafHom{\mathop{\mathcal{H}\!\mathit{om}}\nolimits}
\def\SheafExt{\mathop{\mathcal{E}\!\mathit{xt}}\nolimits}
\def\Sch{\mathit{Sch}}
\def\Mor{\mathop{\mathrm{Mor}}\nolimits}
\def\Ob{\mathop{\mathrm{Ob}}\nolimits}
\def\Sh{\mathop{\mathit{Sh}}\nolimits}
\def\NL{\mathop{N\!L}\nolimits}
\def\CH{\mathop{\mathrm{CH}}\nolimits}
\def\proetale{{pro\text{-}\acute{e}tale}}
\def\etale{{\acute{e}tale}}
\def\QCoh{\mathit{QCoh}}
\def\Ker{\mathop{\mathrm{Ker}}}
\def\Im{\mathop{\mathrm{Im}}}
\def\Coker{\mathop{\mathrm{Coker}}}
\def\Coim{\mathop{\mathrm{Coim}}}

% Boxtimes
%
\DeclareMathSymbol{\boxtimes}{\mathbin}{AMSa}{"02}

%
% Macros for moduli stacks/spaces
%
\def\QCohstack{\mathcal{QC}\!\mathit{oh}}
\def\Cohstack{\mathcal{C}\!\mathit{oh}}
\def\Spacesstack{\mathcal{S}\!\mathit{paces}}
\def\Quotfunctor{\mathrm{Quot}}
\def\Hilbfunctor{\mathrm{Hilb}}
\def\Curvesstack{\mathcal{C}\!\mathit{urves}}
\def\Polarizedstack{\mathcal{P}\!\mathit{olarized}}
\def\Complexesstack{\mathcal{C}\!\mathit{omplexes}}
% \Pic is the operator that assigns to X its picard group, usage \Pic(X)
% \Picardstack_{X/B} denotes the Picard stack of X over B
% \Picardfunctor_{X/B} denotes the Picard functor of X over B
\def\Pic{\mathop{\mathrm{Pic}}\nolimits}
\def\Picardstack{\mathcal{P}\!\mathit{ic}}
\def\Picardfunctor{\mathrm{Pic}}
\def\Deformationcategory{\mathcal{D}\!\mathit{ef}}
\usepackage{cleveref}
\crefname{originalproposition}{original proposition}{original propositions}
\Crefname{originalproposition}{Original proposition}{Original propositions}
\makeatletter
\newlabel{cohomology-proposition-vanishing-Noetherian}{{02UZ}{source}{}{sourcepointer-1}{}}
\newlabel{cohomology-proposition-vanishing-Noetherian@cref}{{[originalproposition][0][]02UZ}{[1][0][]source}{}{}{}}
\makeatother
\crefname{originalequation}{original equation}{original equations}
\Crefname{originalequation}{Original equation}{Original equations}
\makeatletter
\newlabel{equation-wedge}{{0FM2}{source}{}{sourcepointer-2}{}}
\newlabel{equation-wedge@cref}{{[originalequation][0][]0FM2}{[1][0][]source}{}{}{}}
\makeatother
\crefname{originallemma}{original lemma}{original lemmas}
\Crefname{originallemma}{Original lemma}{Original lemmas}
\makeatletter
\newlabel{lemma-bottom-part-degenerates}{{0FW5}{source}{}{sourcepointer-3}{}}
\newlabel{lemma-bottom-part-degenerates@cref}{{[originallemma][0][]0FW5}{[1][0][]source}{}{}{}}
\makeatother
\crefname{originallemma}{original lemma}{original lemmas}
\Crefname{originallemma}{Original lemma}{Original lemmas}
\makeatletter
\newlabel{lemma-duality-hodge}{{0FW4}{source}{}{sourcepointer-4}{}}
\newlabel{lemma-duality-hodge@cref}{{[originallemma][0][]0FW4}{[1][0][]source}{}{}{}}
\makeatother
\crefname{originallemma}{original lemma}{original lemmas}
\Crefname{originallemma}{Original lemma}{Original lemmas}
\makeatletter
\newlabel{lemma-top-part-degenerates}{{0FW6}{source}{}{sourcepointer-5}{}}
\newlabel{lemma-top-part-degenerates@cref}{{[originallemma][0][]0FW6}{[1][0][]source}{}{}{}}
\makeatother
\crefname{originalsection}{original section}{original sections}
\Crefname{originalsection}{Original section}{Original sections}
\makeatletter
\newlabel{section-cup-product}{{0FM1}{source}{}{sourcepointer-6}{}}
\newlabel{section-cup-product@cref}{{[originalsection][0][]0FM1}{[1][0][]source}{}{}{}}
\makeatother
\crefname{originalsection}{original section}{original sections}
\Crefname{originalsection}{Original section}{Original sections}
\makeatletter
\newlabel{section-hodge-to-de-rham}{{0FM6}{source}{}{sourcepointer-7}{}}
\newlabel{section-hodge-to-de-rham@cref}{{[originalsection][0][]0FM6}{[1][0][]source}{}{}{}}
\makeatother
\title{Poincare duality for the de Rham cohomology of a smooth proper equidimensional scheme over a field}
\author{The Stacks Project Authors}
\date{stacks-project:derham.tex:local-snapshot-20261012; GNU FDL 1.2 or later}
\begin{document}
\maketitle
\noindent\textit{Editorial note.} All delivered mathematical text is the author's own, copied from the pinned Stacks Project chapter file derham.tex; no statement or proof step was rewritten, reordered or removed.\par\medskip
\noindent\textit{Editorial note.} Preamble provenance: the pinned chapter file opens with \textbackslash{}input\{preamble\} and carries no preamble text of its own, so the project's single Stacks-wide preamble staged at sources/stacks-derham/preamble.tex is excerpted verbatim as three spans, lines 8-12 (verbatim package and the project's own reference/slogan/history display environments), lines 14-18 (Xy-pic setup) and lines 154-217 (theorem declarations and mathematical macros). The wrapper class and the amsmath/amssymb/amsthm/graphicx base packages are supplied by the assembler. No macro was invented or redefined.\par\medskip
\noindent\textit{Editorial note.} Same-chapter pointers are located at their \textbackslash{}label lines inside the pinned chapter file; cross-chapter pointers are located in the project's own tag list staged in the snapshot (sources/stacks-chow/tags.txt), at the verbatim lines carrying their canonical tags.\par\medskip
\noindent\textit{Allowed standard prerequisites (editorial).}\begin{itemize}
\item The Hodge-to-de Rham spectral sequence. Same chapter derham.tex, Section \textbackslash{}ref\{section-hodge-to-de-rham\} (invoked at proof line 4709). It relates the de Rham cohomology of the smooth proper scheme to the cohomology of the sheaves of differentials, which is how the extreme degrees are identified.
\item Vanishing of the cohomology of a coherent sheaf in degrees above the dimension, and the degeneration of the extreme parts of the spectral sequence. Cohomology, label cohomology-proposition-vanishing-Noetherian (invoked at proof line 4711) and Same chapter derham.tex, labels lemma-bottom-part-degenerates and lemma-top-part-degenerates (invoked at proof lines 4712 and 4713). Together with the vanishing of \$\textbackslash{}Omega\textasciicircum{}i\$ for \$i > d\$ they give \$H\textasciicircum{}0\_\{dR\}(X/k) = H\textasciicircum{}0(X, \textbackslash{}mathcal\{O\}\_X)\$ and \$H\textasciicircum{}d(X, \textbackslash{}Omega\textasciicircum{}d\_\{X/k\}) = H\textasciicircum{}\{2d\}\_\{dR\}(X/k)\$.
\item The Hodge duality pairing in the extreme case. Same chapter derham.tex, label lemma-duality-hodge (invoked at proof lines 4724 and 4768). It supplies the trace map in the base case \$p = 0\$ of the induction, giving perfect pairings between \$H\textasciicircum{}i(X, \textbackslash{}mathcal\{O\}\_X)\$ and \$H\textasciicircum{}\{d-i\}(X, \textbackslash{}Omega\textasciicircum{}d)\$.
\item The wedge product on the total complex and the construction of cup products. Same chapter derham.tex, equation equation-wedge (invoked at proof line 4731) and Section \textbackslash{}ref\{section-cup-product\} (invoked at proof line 4752). The wedge map on \$\textbackslash{}operatorname\{Tot\}(\textbackslash{}Omega\textasciicircum{}\textbackslash{}bullet \textbackslash{}otimes \textbackslash{}Omega\textasciicircum{}\textbackslash{}bullet)\$ restricts to the maps \$\textbackslash{}gamma\_p\$ whose induced products are the pairings studied in the induction.
\end{itemize}\medskip
\ifdefined\hypertarget\hypertarget{sourcepointer-1}{}\fi
\noindent\textit{Source reference.} Original proposition 02UZ: sources/stacks-chow/tags.txt, source lines 3750--3750. Navigation pointer for \textbackslash{}ref\{cohomology-proposition-vanishing-Noetherian\} used inside the counted span (line 4711); the argument that consumes it is retained in the span.\par\medskip
\ifdefined\hypertarget\hypertarget{sourcepointer-2}{}\fi
\noindent\textit{Source reference.} Original equation 0FM2: sources/stacks-derham/pinned-original-source.tex, source lines 495--495. Navigation pointer for \textbackslash{}ref\{equation-wedge\} used inside the counted span (line 4731); the argument that consumes it is retained in the span.\par\medskip
\ifdefined\hypertarget\hypertarget{sourcepointer-3}{}\fi
\noindent\textit{Source reference.} Original lemma 0FW5: sources/stacks-derham/pinned-original-source.tex, source lines 4548--4548. Navigation pointer for \textbackslash{}ref\{lemma-bottom-part-degenerates\} used inside the counted span (line 4712); the argument that consumes it is retained in the span.\par\medskip
\ifdefined\hypertarget\hypertarget{sourcepointer-4}{}\fi
\noindent\textit{Source reference.} Original lemma 0FW4: sources/stacks-derham/pinned-original-source.tex, source lines 4496--4496. Navigation pointer for \textbackslash{}ref\{lemma-duality-hodge\} used inside the counted span (line 4724); the argument that consumes it is retained in the span.\par\medskip
\ifdefined\hypertarget\hypertarget{sourcepointer-5}{}\fi
\noindent\textit{Source reference.} Original lemma 0FW6: sources/stacks-derham/pinned-original-source.tex, source lines 4577--4577. Navigation pointer for \textbackslash{}ref\{lemma-top-part-degenerates\} used inside the counted span (line 4713); the argument that consumes it is retained in the span.\par\medskip
\ifdefined\hypertarget\hypertarget{sourcepointer-6}{}\fi
\noindent\textit{Source reference.} Original section 0FM1: sources/stacks-derham/pinned-original-source.tex, source lines 483--483. Navigation pointer for \textbackslash{}ref\{section-cup-product\} used inside the counted span (line 4752); the argument that consumes it is retained in the span.\par\medskip
\ifdefined\hypertarget\hypertarget{sourcepointer-7}{}\fi
\noindent\textit{Source reference.} Original section 0FM6: sources/stacks-derham/pinned-original-source.tex, source lines 671--671. Navigation pointer for \textbackslash{}ref\{section-hodge-to-de-rham\} used inside the counted span (line 4709); the argument that consumes it is retained in the span.\par\medskip
\begin{proposition}
\label{proposition-poincare-duality}
Let $k$ be a field. Let $X$ be a nonempty smooth proper scheme over $k$
equidimensional of dimension $d$. There exists a $k$-linear map
$$
t : H^{2d}_{dR}(X/k) \longrightarrow k
$$
unique up to precomposing by multiplication by a unit of
$H^0(X, \mathcal{O}_X)$ with the following property: for all $i$ the pairing
$$
H^i_{dR}(X/k) \times H_{dR}^{2d - i}(X/k)
\longrightarrow
k, \quad
(\xi, \xi') \longmapsto t(\xi \cup \xi')
$$
is perfect.
\end{proposition}

\begin{proof}
By the Hodge-to-de Rham spectral sequence
(Section \ref{section-hodge-to-de-rham}), the vanishing
of $\Omega^i_{X/k}$ for $i > d$, the vanishing in
Cohomology, Proposition \ref{cohomology-proposition-vanishing-Noetherian}
and the results of Lemmas \ref{lemma-bottom-part-degenerates} and
\ref{lemma-top-part-degenerates}
we see that $H^0_{dR}(X/k) = H^0(X, \mathcal{O}_X)$
and $H^d(X, \Omega^d_{X/k}) = H_{dR}^{2d}(X/k)$.
More precisely, these identifications come from the maps
of complexes
$$
\Omega^\bullet_{X/k} \to \mathcal{O}_X[0]
\quad\text{and}\quad
\Omega^d_{X/k}[-d] \to \Omega^\bullet_{X/k}
$$
Let us choose $t : H_{dR}^{2d}(X/k) \to k$ which via this identification
corresponds to a $t$ as in Lemma \ref{lemma-duality-hodge}.
Then in any case we see that the pairing displayed in the lemma
is perfect for $i = 0$.

\medskip\noindent
Denote $\underline{k}$ the constant sheaf with value $k$ on $X$.
Let us abbreviate $\Omega^\bullet = \Omega^\bullet_{X/k}$.
Consider the map (\ref{equation-wedge}) which in our situation reads
$$
\wedge :
\text{Tot}(\Omega^\bullet \otimes_{\underline{k}} \Omega^\bullet)
\longrightarrow
\Omega^\bullet
$$
For every integer $p = 0, 1, \ldots, d$ this map
annihilates the subcomplex
$\text{Tot}(\sigma_{> p} \Omega^\bullet \otimes_{\underline{k}}
\sigma_{\geq d - p} \Omega^\bullet)$ for degree reasons.
Hence we find that the restriction of $\wedge$ to the subcomplex
$\text{Tot}(\Omega^\bullet \otimes_{\underline{k}}
\sigma_{\geq d - p}\Omega^\bullet)$ factors through a map of complexes
$$
\gamma_p :
\text{Tot}(\sigma_{\leq p} \Omega^\bullet \otimes_{\underline{k}}
\sigma_{\geq d - p} \Omega^\bullet)
\longrightarrow
\Omega^\bullet
$$
Using the same procedure as in Section \ref{section-cup-product} we obtain
cup products
$$
H^i(X, \sigma_{\leq p} \Omega^\bullet) \times
H^{2d - i}(X, \sigma_{\geq d - p}\Omega^\bullet)
\longrightarrow
H_{dR}^{2d}(X, \Omega^\bullet)
$$
We will prove by induction on $p$ that these cup products via $t$
induce perfect pairings between $H^i(X, \sigma_{\leq p} \Omega^\bullet)$
and $H^{2d - i}(X, \sigma_{\geq d - p}\Omega^\bullet)$. For $p = d$
this is the assertion of the proposition.

\medskip\noindent
The base case is $p = 0$. In this case we simply obtain the pairing
between $H^i(X, \mathcal{O}_X)$ and $H^{d - i}(X, \Omega^d)$ of
Lemma \ref{lemma-duality-hodge} and the result is true.

\medskip\noindent
Induction step. Say we know the result is true for $p$. Then
we consider the distinguished triangle
$$
\Omega^{p + 1}[-p - 1] \to
\sigma_{\leq p + 1}\Omega^\bullet \to
\sigma_{\leq p}\Omega^\bullet \to
\Omega^{p + 1}[-p]
$$
and the distinguished triangle
$$
\sigma_{\geq d - p}\Omega^\bullet \to
\sigma_{\geq d - p - 1}\Omega^\bullet \to
\Omega^{d - p - 1}[-d + p + 1] \to
(\sigma_{\geq d - p}\Omega^\bullet)[1]
$$
Observe that both are distinguished triangles in the homotopy category
of complexes of sheaves of $\underline{k}$-modules; in particular the
maps $\sigma_{\leq p}\Omega^\bullet \to \Omega^{p + 1}[-p]$ and
$\Omega^{d - p - 1}[-d + p + 1] \to (\sigma_{\geq d - p}\Omega^\bullet)[1]$
are given by actual maps of complexes, namely using the differential
$\Omega^p \to \Omega^{p + 1}$ and the differential
$\Omega^{d - p - 1} \to \Omega^{d - p}$.
Consider the long exact cohomology sequences associated to these
distinguished triangles
$$
\xymatrix{
H^{i - 1}(X, \sigma_{\leq p}\Omega^\bullet) \ar[d]_a \\
H^i(X, \Omega^{p + 1}[-p - 1]) \ar[d]_b \\
H^i(X, \sigma_{\leq p + 1}\Omega^\bullet) \ar[d]_c \\
H^i(X, \sigma_{\leq p}\Omega^\bullet) \ar[d]_d \\
H^{i + 1}(X, \Omega^{p + 1}[-p - 1])
}
\quad\quad
\xymatrix{
H^{2d - i  + 1}(X, \sigma_{\geq d - p}\Omega^\bullet) \\
H^{2d - i}(X, \Omega^{d - p - 1}[-d + p + 1]) \ar[u]_{a'} \\
H^{2d - i}(X, \sigma_{\geq d - p - 1}\Omega^\bullet) \ar[u]_{b'} \\
H^{2d - i}(X, \sigma_{\geq d - p}\Omega^\bullet) \ar[u]_{c'} \\
H^{2d - i - 1}(X, \Omega^{d - p - 1}[-d + p + 1]) \ar[u]_{d'}
}
$$
By induction and Lemma \ref{lemma-duality-hodge}
we know that the pairings constructed above between the
$k$-vectorspaces on the first, second, fourth, and fifth
rows are perfect. By the $5$-lemma, in order to show that
the pairing between the cohomology groups in the middle row
is perfect, it suffices to show that the pairs
$(a, a')$, $(b, b')$, $(c, c')$, and $(d, d')$
are compatible with the given pairings (see below).

\medskip\noindent
Let us prove this for the pair $(c, c')$. Here we observe simply
that we have a commutative diagram
$$
\xymatrix{
\text{Tot}(\sigma_{\leq p} \Omega^\bullet \otimes_{\underline{k}}
\sigma_{\geq d - p} \Omega^\bullet) \ar[d]_{\gamma_p} &
\text{Tot}(\sigma_{\leq p + 1} \Omega^\bullet \otimes_{\underline{k}}
\sigma_{\geq d - p} \Omega^\bullet) \ar[l] \ar[d] \\
\Omega^\bullet &
\text{Tot}(\sigma_{\leq p + 1} \Omega^\bullet \otimes_{\underline{k}}
\sigma_{\geq d - p - 1} \Omega^\bullet) \ar[l]_-{\gamma_{p + 1}}
}
$$
Hence if we have $\alpha \in H^i(X, \sigma_{\leq p + 1}\Omega^\bullet)$
and $\beta \in H^{2d - i}(X, \sigma_{\geq d - p}\Omega^\bullet)$
then we get
$\gamma_p(\alpha \cup c'(\beta)) = \gamma_{p + 1}(c(\alpha) \cup \beta)$
by functoriality of the cup product.

\medskip\noindent
Similarly for the pair $(b, b')$ we use the commutative diagram
$$
\xymatrix{
\text{Tot}(\sigma_{\leq p + 1} \Omega^\bullet \otimes_{\underline{k}}
\sigma_{\geq d - p - 1} \Omega^\bullet) \ar[d]_{\gamma_{p + 1}} &
\text{Tot}(\Omega^{p + 1}[-p - 1] \otimes_{\underline{k}}
\sigma_{\geq d - p - 1} \Omega^\bullet) \ar[l] \ar[d] \\
\Omega^\bullet &
\Omega^{p + 1}[-p - 1]
\otimes_{\underline{k}}
\Omega^{d - p - 1}[-d + p + 1] \ar[l]_-\wedge
}
$$
and argue in the same manner.

\medskip\noindent
For the pair $(d, d')$ we use the commutative diagram
$$
\xymatrix{
\Omega^{p + 1}[-p] \otimes_{\underline{k}}
\Omega^{d - p - 1}[-d + p] \ar[d] &
\text{Tot}(\sigma_{\leq p}\Omega^\bullet \otimes_{\underline{k}}
\Omega^{d - p - 1}[-d + p]) \ar[l] \ar[d] \\
\Omega^\bullet &
\text{Tot}(\sigma_{\leq p}\Omega^\bullet \otimes_{\underline{k}}
\sigma_{\geq d - p}\Omega^\bullet) \ar[l]
}
$$
and we look at cohomology classes in
$H^i(X, \sigma_{\leq p}\Omega^\bullet)$ and
$H^{2d - i}(X, \Omega^{d - p - 1}[-d + p])$.
Changing $i$ to $i - 1$ we get the result for the pair $(a, a')$
thereby finishing the proof that our pairings are perfect.

\medskip\noindent
We omit the argument showing the uniqueness of $t$ up to
precomposing by multiplication by a unit in $H^0(X, \mathcal{O}_X)$.
\end{proof}

\end{document}
